\section*{Chapitre \ref{ch:b2:heart} --- Le cœur et la révolution cardiaque}

\begin{solution}{exo:b2:heart:1}
\omterm{def:b2:blood-circulation:vessels}{Veine} cave $\to$ \omterm{def:b2:heart:anatomy}{oreillette} droite $\to$ valve tricuspide $\to$ \omterm{def:b2:heart:anatomy}{ventricule}
droit $\to$ valve pulmonaire $\to$ \omterm{def:b2:blood-circulation:vessels}{artère} pulmonaire $\to$ poumons
$\to$ \omterm{def:b2:blood-circulation:vessels}{veines} pulmonaires $\to$ \omterm{def:b2:heart:anatomy}{oreillette} gauche $\to$ valve mitrale
$\to$ \omterm{def:b2:heart:anatomy}{ventricule} gauche $\to$ valve aortique $\to$ aorte.
\end{solution}

\begin{solution}{exo:b2:heart:2}
Remplissage (diastole)~: valves auriculo-ventriculaires ouvertes, valves
sigmoïdes fermées. Contraction isovolumétrique~: les quatre valves
fermées. Éjection~: valves sigmoïdes ouvertes, auriculo-ventriculaires
fermées. Relaxation isovolumétrique~: toutes fermées. Premier bruit~: la
fermeture des valves auriculo-ventriculaires au début de la systole~;
second bruit~: la fermeture des valves sigmoïdes à sa fin.
\end{solution}

\begin{solution}{exo:b2:heart:3}
\omterm{prop:b2:heart:conduction}{Nœud sinusal} ($t = 0$~; l'onde P quand les \omterm{def:b2:heart:anatomy}{oreillettes} se
dépolarisent)~; \omterm{prop:b2:heart:conduction}{nœud atrioventriculaire} (atteint vers \qty{80}{ms}, qui
retient l'influx pendant \qty{100}{ms}~: le segment PR)~; faisceau de His
et \omterm{prop:b2:heart:conduction}{fibres de Purkinje} (\qtyrange{180}{220}{ms}~; le QRS quand les
\omterm{def:b2:heart:anatomy}{ventricules} se dépolarisent)~; l'onde T plus tard, lorsqu'ils se
repolarisent.
\end{solution}

\begin{solution}{exo:b2:heart:4}
Dans sa plage de fonctionnement, le \omterm{def:b2:heart:anatomy}{ventricule} éjecte davantage lorsqu'il
a été davantage rempli. Si le \omterm{def:b2:heart:anatomy}{ventricule} droit pompe davantage, le gauche
reçoit davantage quelques battements plus tard et, en vertu de la même
loi, pompe davantage~: tout déséquilibre se corrige de lui-même, sans
signal.
\end{solution}

\begin{solution}{exo:b2:heart:5}
Volume d'éjection \qty{80}{mL}~; \omterm{prop:b2:heart:cycle}{fraction d'éjection}
$80/130 = \qty{62}{\%}$~; débit $80\times 70 = \qty{5.6}{L/min}$. \omterm{def:b2:heart:anatomy}{Cœur}
défaillant~: \qty{60}{mL}, $60/180 = \qty{33}{\%}$, \qty{4.2}{L/min} ---
le \omterm{def:b2:heart:anatomy}{ventricule} dilaté maintient à peu près le débit en travaillant à grand
volume, mais n'éjecte que le tiers de ce qu'il contient.
\end{solution}

\begin{solution}{exo:b2:heart:6}
$W = 100\times 133\times 7\times 10^{-5} = \qty{0.93}{J}$~; puissance
$0.93\times 1.2 = \qty{1.1}{W}$. À \qty{150}{mmHg}~: \qty{1.4}{J},
\qty{1.7}{W}. Les \qty{0.56}{W} mécaniques supplémentaires représentent
\qty{3.7}{W} métaboliques, soit \qty{0.19}{mL} de dioxygène par
seconde~: \qty{11}{mL/min} de plus.
\end{solution}

\begin{solution}{exo:b2:heart:7}
$60/0.5 = 120$ battements par minute~; $60/1.5 = 40$. Troisième cas~: un
bloc auriculo-ventriculaire complet --- les \omterm{def:b2:heart:anatomy}{oreillettes} battent à 100 sous
la commande du \omterm{prop:b2:heart:conduction}{nœud sinusal}, les \omterm{def:b2:heart:anatomy}{ventricules} à 33 sous celle d'un
pacemaker qui leur est propre, et le \omterm{prop:b2:heart:conduction}{nœud atrioventriculaire} ne conduit
plus rien.
\end{solution}

\begin{solution}{exo:b2:heart:8}
Dérive~: $20/25 = \qty{0.8}{s}$, plus \qty{0.15}{s}~: \qty{0.95}{s}, soit
63 battements par minute. Acétylcholine~: $25/18 = \qty{1.39}{s}$, plus
$0.15$~: \qty{1.54}{s}, 39 par minute. Noradrénaline~:
$20/40 = \qty{0.5}{s}$, plus $0.15$~: \qty{0.65}{s}, 92 par minute.
\end{solution}

\begin{solution}{exo:b2:heart:9}
Le retard laisse aux \omterm{def:b2:heart:anatomy}{oreillettes} le temps d'achever de se vider dans les
\omterm{def:b2:heart:anatomy}{ventricules} avant que ceux-ci ne se contractent~; sans lui, \omterm{def:b2:heart:anatomy}{oreillettes}
et \omterm{def:b2:heart:anatomy}{ventricules} se contracteraient ensemble, les valves
auriculo-ventriculaires se fermeraient sur des \omterm{def:b2:heart:anatomy}{ventricules} à moitié
remplis et le volume d'éjection perdrait la contribution des
\omterm{def:b2:heart:anatomy}{oreillettes}. Un long intervalle PR ne fait que retarder la systole
ventriculaire et ne coûte rien, jusqu'à ce que le retard devienne si
long que des battements sautent ou que la conduction échoue
complètement.
\end{solution}

\begin{solution}{exo:b2:heart:10}
Repos~: $5000/45 = \qty{111}{mL}$~; effort~: $\num{30000}/180 =
\qty{167}{mL}$. La loi de Starling convertit le retour veineux accru de
l'effort en un volume d'éjection plus grand~; les nerfs sympathiques
ajoutent la fréquence et la contractilité (un volume télésystolique plus
faible). Au-delà de 200 environ, la diastole est trop courte pour remplir
le \omterm{def:b2:heart:anatomy}{ventricule}, et le volume d'éjection diminue plus vite que la fréquence
n'augmente.
\end{solution}

\begin{solution}{exo:b2:heart:11}
$v = Q/A$~: $250/3 = \qty{83}{cm/s}$ et $250/1 = \qty{250}{cm/s}$.
$\Delta P = \tfrac12\times 1060\times v^{2}$~: \qty{365}{Pa}
(\qty{2.7}{mmHg}) et \qty{3300}{Pa} (\qty{25}{mmHg}). Le \omterm{def:b2:heart:anatomy}{ventricule}
produit la pression supplémentaire en épaississant sa paroi~; ce
\omterm{def:b2:heart:anatomy}{ventricule} épais et rigide finit par dépasser les capacités de son
irrigation coronaire, se remplit mal et défaille.
\end{solution}

\begin{solution}{exo:b2:heart:12}
L'automatisme fournit un battement sans aucune commande~; la loi de
Starling ajuste le débit au retour veineux et équilibre les deux
\omterm{def:b2:heart:anatomy}{ventricules} sans aucune commande~; les nerfs et l'adrénaline ne font que
régler la fréquence et la contractilité autour de ce noyau autorégulé.
Un \omterm{def:b2:heart:anatomy}{cœur} transplanté bat (à environ 100, sans tonus vagal), ajuste son
volume d'éjection à son remplissage, et peut augmenter son débit à
l'effort par le mécanisme de Starling et par l'adrénaline circulante ---
mais lentement, et moins qu'un \omterm{def:b2:heart:anatomy}{cœur} innervé.
\end{solution}

\begin{solution}{pb:b2:heart:1}
\textbf{1.} \qty{70}{mL}~; $70/130 = \qty{54}{\%}$~; $70\times 70 =
\qty{4.9}{L/min}$.
\textbf{2.} $60/70 = \qty{0.86}{s}$~; diastole $0.86 - 0.3 =
\qty{0.56}{s}$.
\textbf{3.} $70/0.3 = \qty{233}{mL/s}$~; $233/3 = \qty{78}{cm/s}$.
\textbf{4.} $70/1500 = \qty{0.047}{s}$, environ \qty{50}{ms}.
\textbf{5.} $100\times 133\times 7\times 10^{-5} = \qty{0.93}{J}$~;
$0.93\times 70/60 = \qty{1.1}{W}$.
\textbf{6.} Les deux \omterm{def:b2:heart:anatomy}{ventricules}~: $1.1\times 1.2 = \qty{1.3}{W}$~; à
\qty{15}{\%}, \qty{8.7}{W} métaboliques~; $8.7/20 = \qty{0.43}{mL}$ de
dioxygène par seconde, \qty{26}{mL/min}.
\textbf{7.} Avec une extraction de \qty{0.14}{mL} par millilitre de
\omterm{def:b2:blood-circulation:blood}{sang}~: $26/0.14 = \qty{190}{mL/min}$.
\textbf{8.} Intervalle \qty{0.86}{s}, durée de dérive \qty{0.71}{s}~:
$20/0.71 = \qty{28}{mV/s}$.
\textbf{9.} Intervalle \qty{0.6}{s}, dérive \qty{0.45}{s}~:
\qty{44}{mV/s}.
\textbf{10.} Intervalle \qty{0.33}{s}, dérive \qty{0.18}{s}~:
\qty{110}{mV/s}.
\textbf{11.} R--R \qty{0.86}{s}~; P~: dépolarisation des \omterm{def:b2:heart:anatomy}{oreillettes}~;
QRS~: dépolarisation ventriculaire~; T~: repolarisation ventriculaire.
\textbf{12.} Le retard laisse les \omterm{def:b2:heart:anatomy}{oreillettes} se vider dans les
\omterm{def:b2:heart:anatomy}{ventricules} avant leur contraction~; à 180, le cycle entier dure
\qty{0.33}{s}, si bien qu'un retard inchangé de \qty{0.16}{s} en
occuperait la moitié --- les nerfs sympathiques accélèrent aussi la
conduction du nœud.
\textbf{13.} Bloc auriculo-ventriculaire complet~: les \omterm{def:b2:heart:anatomy}{ventricules} à 40
par minute, rythmés par le faisceau ou les \omterm{prop:b2:heart:conduction}{fibres de Purkinje}, les
\omterm{def:b2:heart:anatomy}{oreillettes} à 75 sous la commande du \omterm{prop:b2:heart:conduction}{nœud sinusal}.
\textbf{14.} \qty{120}{mL}~; $120/150 = \qty{80}{\%}$~; $120\times 180 =
\qty{21.6}{L/min}$.
\textbf{15.} $0.33 - 0.3 = \qty{0.03}{s}$~: presque pas de temps pour se
remplir~; plus vite encore, le volume d'éjection s'effondre.
\textbf{16.} $W = 120\times 133\times 1.2\times 10^{-4} = \qty{1.9}{J}$,
$\times 3$ par seconde~: \qty{5.7}{W}~; les deux \omterm{def:b2:heart:anatomy}{ventricules}
\qty{6.9}{W}~; métabolisme \qty{46}{W}~; dioxygène \qty{2.3}{mL/s}, soit
\qty{140}{mL/min}.
\textbf{17.} $140/0.14 = \qty{1000}{mL/min}$, cinq fois le débit de
repos, à délivrer pendant une diastole réduite au dixième du cycle~: les
\omterm{def:b2:blood-circulation:vessels}{artérioles} coronaires doivent se dilater au maximum.
\textbf{18.} Fréquence d'environ 100 sans tonus vagal, volume
télésystolique inchangé à \qty{60}{mL}~: volume d'éjection
$150 - 60 = \qty{90}{mL}$, débit \qty{9}{L/min}.
\textbf{19.} Fréquence seule ($70 \to 180$ à \qty{70}{mL}, soit
\qty{12.6}{L/min})~: $+7.7$~; remplissage (\qty{20}{mL} de plus à
180)~: $+3.6$~; contractilité (\qty{30}{mL} de moins de volume résiduel
à 180)~: $+5.4$ --- de 4.9 à \qty{21.6}{L/min}, les trois contributions
totalisant exactement $+16.7$.
\textbf{20.} $233/1 = \qty{233}{cm/s}$, soit \qty{2.3}{m/s}.
\textbf{21.} $\tfrac12\times 1060\times 2.33^{2} = \qty{2900}{Pa}$, soit
\qty{22}{mmHg}.
\textbf{22.} $120/0.3 = \qty{400}{mL/s}$, \qty{4}{m/s}~; $\tfrac12\times
1060\times 16 = \qty{8500}{Pa}$, soit \qty{64}{mmHg}.
\textbf{23.} Pression maximale $120 + 22 = \qty{142}{mmHg}$ au repos,
environ $140 +
64 = \qty{204}{mmHg}$ à l'effort~; la pression moyenne d'éjection au
repos passe de 100 à 122~: travail par battement $\times 1.22$.
\textbf{24.} Une paroi plus épaisse développe plus de force et réduit la
contrainte par fibre~; mais la masse à irriguer augmente, tandis que le
débit coronaire, écrasé par la paroi épaisse en systole et disposant
d'une diastole plus courte, ne suit pas, et une paroi épaisse se relâche
mal et se remplit moins~: l'ischémie et la défaillance s'ensuivent.
\textbf{25.} \qty{4.9}{L/min} au repos, \qty{21.6}{L/min} à l'effort~;
\qty{0.93}{J} par battement~; gradient de \qty{22}{mmHg} au repos et de
\qty{64}{mmHg} à l'effort.
\end{solution}
